\section{Análises Finais}

Este estudo buscou quantificar o impacto da escolha da linguagem de programação no desempenho de algoritmos de visão computacional. As evidências coletadas permitem concluir que o ``custo da abstração'' é real e mensurável, representando um acréscimo de latência entre 13\% e 20\% para linguagens gerenciadas (Java e Python) em comparação ao C++ nativo nos cenários testados.

A análise comparativa permite traçar as seguintes diretrizes para a engenharia de \textit{software} em visão computacional:

\begin{itemize}
    \item \textbf{C++:} Permanece a escolha ideal para sistemas embarcados, aplicações críticas de segurança e cenários onde o \textit{throughput} e o determinismo temporal são prioritários, justificando sua curva de aprendizado mais íngreme.
    
    \item \textbf{Python:} Consolida-se como a melhor ferramenta para pesquisa, desenvolvimento rápido e aplicações em servidor onde a escalabilidade horizontal é viável. A percepção positiva da comunidade e a vastidão do ecossistema compensam a latência adicional em cenários não-críticos.
    
    \item \textbf{Java:} Apresenta-se como uma opção não ideal para aplicações de visão computacional devido à complexidade do ambiente de execução, falta de suporte e à imprevisibilidade introduzida pelo \textit{garbage collector}. Seu uso deve ser cuidadosamente avaliado frente aos requisitos de desempenho e estabilidade do sistema.
\end{itemize}

Como trabalhos futuros, sugere-se a extensão da análise para outras bibliotecas populares de visão computacional, como TensorFlow, PyTorch e Keras, bem como a inclusão de métricas adicionais relacionadas ao consumo de memória, eficiência energética e escalabilidade em ambientes distribuídos. Ademais, investigações sobre o impacto de diferentes paradigmas de programação (funcional, orientada a objetos, etc.) no desempenho de algoritmos de visão computacional podem oferecer insights valiosos para a comunidade.

Em suma, não existe uma "melhor linguagem" universal para o OpenCV; a escolha deve ser pautada pelo equilíbrio entre os requisitos não-funcionais de desempenho do sistema e os requisitos de negócio relacionados ao tempo de entrega e manutenibilidade do código.